\HandoutSetup
  {NE 630}
  {03}
  {Radioactive Decay}
  {Read FNRP Sections 1.6--1.7 and define the key terms before class.}

\begin{document}
\MakeHandoutTitle

\begin{handoutcolumns}

\begin{fixedbox}{Macro objective}
Determine the amounts of a radioactive species and its daughters at any point
in time from initial inventories, production rates, and decay constants.
\end{fixedbox}

\TermStrip{radioactivity; $\lambda$; half-life; mean lifetime; activity; Bq; Ci; initial-value problem; saturation; steady state; secular equilibrium}

\HandoutSection{1}{Why radioactive species appear}

\begin{fixedbox}{Fission-product bridge from Lesson 2}
Fission fragments are neutron rich. Their $\beta^-$ decays produce delayed
$\beta$, $\bar\nu$, and $\gamma$ energy. When a daughter state is formed above
its neutron-separation energy, one decay branch can emit a delayed neutron:
\scriptsize
\begin{align*}
{}^{87}_{35}\mathrm{Br}
 &\xrightarrow[55.65\unit{s}]{\beta^-}
 {}^{87}_{36}\mathrm{Kr}^{*}
 \xrightarrow{n}{}^{86}_{36}\mathrm{Kr},\\[-1pt]
{}^{137}_{53}\mathrm{I}
 &\xrightarrow[24.5\unit{s}]{\beta^-}
 {}^{137}_{54}\mathrm{Xe}^{*}
 \xrightarrow{n}{}^{136}_{54}\mathrm{Xe}.
\end{align*}
\normalsize
The notes give delayed-neutron energies of roughly
\RevealBlank[0.80in]{$0.2$--$0.6\unit{MeV}$}. Neutron emission is a branch,
not the result of every $\beta^-$ decay. Deposited $\beta$ and $\gamma$ energy
drives spent-fuel decay heat; antineutrino energy escapes.
\end{fixedbox}

\HandoutSection{2}{The single-nuclide decay law}

\begin{boardbox}{From a short-time probability to $N(t)$}
For $N(t)$ radioactive nuclei,
\[
 \Delta N\approx-\lambda N(t)\Delta t
 \quad\Longrightarrow\quad
 \boxed{N'=-\lambda N},\qquad N(0)=N_0.
\]
Separate and integrate:
\[
 \frac{dN}{N}=-\lambda\,dt,
 \qquad
 \ln\!\frac{N(t)}{N_0}=\Blank[0.72in],
\]
\[
 \boxed{N(t)=\MathBlank[1.25in]{N_0e^{-\lambda t}}}.
\]
Here $\lambda$ is the decay probability per unit time for one nucleus.
\end{boardbox}

\begin{checkpoint}[Half-life check]
\[
 \frac{N(nt_{1/2})}{N_0}=\MathBlank[0.72in]{2^{-n}},
 \qquad
 \frac{A(nt_{1/2})}{A_0}=\MathBlank[0.72in]{2^{-n}}.
\]
What assumptions make both ratios equal to $2^{-n}$?
\RuledLines{1}
\end{checkpoint}

\switchcolumn

\HandoutSection{3}{Clocks, number, mass, and activity}

\begin{fixedbox}{Fixed relations and unit conversions}
\[
 \boxed{\lambda=\frac{\ln 2}{t_{1/2}}=\frac{1}{\bar t}},
 \qquad
 \boxed{A(t)=\lambda N(t)}.
\]
\scriptsize
\begin{tabularx}{\linewidth}{@{}>{\bfseries}l X@{}}
number--mass & $N=(m/M)N_A$\\
activity & decays per unit time\\
Bq & $1\unit{Bq}=1\unit{s^{-1}}$\\
Ci & $1\unit{Ci}=3.7\times10^{10}\unit{Bq}$
\end{tabularx}
\normalsize
When there is no production, $A(t)=-dN/dt$. Thus, for a pure radionuclide,
\[
 m\ \xrightarrow{\div M}\ \mathrm{mol}
 \ \xrightarrow{\times N_A}\ N
 \ \xrightarrow{\times\lambda}\ A.
\]
\end{fixedbox}

\begin{checkpoint}[Reverse the conversion route]
Express the radionuclide mass in terms of measured activity:
\[
 m=\MathBlank[2.05in]{\frac{AM}{\lambda N_A}}.
\]
List the units you would use for each quantity.
\end{checkpoint}

\HandoutSection{4}{Production and saturation}

\begin{boardbox}{Decay with a source and saturation}
\footnotesize
If nuclei are added at rate $R(t)$,
\[
 \boxed{N'=-\lambda N+R(t)},\qquad N(t_0)\ \text{given}.
\]
The integrating factor $e^{\lambda t}$ gives
\[
 \frac{d}{dt}\!\left(Ne^{\lambda t}\right)=R(t)e^{\lambda t},
\]
\[
 \boxed{\begin{aligned}
 N(t)={}&N(t_0)e^{-\lambda(t-t_0)}\\[-2pt]
 &+\int_{t_0}^{t}R(t')e^{-\lambda(t-t')}\,dt'.
 \end{aligned}}
\]
For $R(t)=R_0$ and $t_0=0$,
\[
 \boxed{N(t)=N_0e^{-\lambda t}
 +\frac{R_0}{\lambda}\left(1-e^{-\lambda t}\right)}.
\]
For $N_0=0$,
\[
 \begin{aligned}
 A(t)&=\MathBlank[1.30in]{R_0(1-e^{-\lambda t})},\\[-2pt]
 N_{\mathrm{sat}}&=\MathBlank[0.60in]{R_0/\lambda},\qquad
 A_{\mathrm{sat}}=\MathBlank[0.46in]{R_0}.
 \end{aligned}
\]
After $n$ half-lives, $A/A_{\mathrm{sat}}=1-2^{-n}$; one, three, five, and
seven half-lives give 50\%, 87.5\%, 96.9\%, and 99.2\% of saturation.
\normalsize
\end{boardbox}

\end{handoutcolumns}

\newpage
\begin{handoutcolumns}

\HandoutSection{5}{The $A=94$ fission-product chain}

\begin{fixedbox}{Chain and initial data from the notes}
\[
{}^{94}_{38}\mathrm{Sr}
 \xrightarrow[75\unit{s}]{\beta^-}
{}^{94}_{39}\mathrm{Y}
 \xrightarrow[18.7\unit{min}]{\beta^-}
{}^{94}_{40}\mathrm{Zr}\quad(\text{stable}).
\]
Assume $N_{\mathrm{Sr}}(0)=N_0$,
$N_{\mathrm{Y}}(0)=N_{\mathrm{Zr}}(0)=0$, and no production after $t=0$.
\end{fixedbox}

\begin{boardbox}{Balance each nuclide}
\begin{align*}
 N_{\mathrm{Sr}}'&=-\lambda_{\mathrm{Sr}}N_{\mathrm{Sr}},\\
 N_{\mathrm{Y}}'&=\lambda_{\mathrm{Sr}}N_{\mathrm{Sr}}
                  -\lambda_{\mathrm{Y}}N_{\mathrm{Y}},\\
 N_{\mathrm{Zr}}'&=\lambda_{\mathrm{Y}}N_{\mathrm{Y}}.
\end{align*}
Label each term as a gain or loss.  Why can this one-way chain be solved from
the parent downstream?
\RuledLines{1}
\end{boardbox}

\begin{fixedbox}{Sequential solution}
\[
 N_{\mathrm{Sr}}(t)=N_0e^{-\lambda_{\mathrm{Sr}}t},
\]
\[
 N_{\mathrm{Y}}(t)
 =\frac{\lambda_{\mathrm{Sr}}N_0}
 {\lambda_{\mathrm{Y}}-\lambda_{\mathrm{Sr}}}
 \left(e^{-\lambda_{\mathrm{Sr}}t}-e^{-\lambda_{\mathrm{Y}}t}\right),
\]
\[
 \begin{aligned}
 N_{\mathrm{Zr}}(t)
 &=\int_0^t\lambda_{\mathrm{Y}}N_{\mathrm{Y}}(t')\,dt'\\[-1pt]
 &=N_0-N_{\mathrm{Sr}}(t)-N_{\mathrm{Y}}(t).
 \end{aligned}
\]
The last line also follows from conservation of the $A=94$ nuclei.
\end{fixedbox}

\begin{checkpoint}[Inventory after one hour]
At $t=1\unit{h}$,
\[
 \frac{N_{\mathrm{Sr}}}{N_0}=\MathBlank[0.52in]{3.6\times10^{-15}},\quad
 \frac{N_{\mathrm{Y}}}{N_0}=\MathBlank[0.52in]{0.116},\quad
 \frac{N_{\mathrm{Zr}}}{N_0}=\MathBlank[0.52in]{0.884}.
\]
Which half-life controls the long-time approach to stable zirconium?
\RuledLines{1}
\end{checkpoint}

\switchcolumn

\HandoutSection{6}{General chains and equilibrium}

\begin{fixedbox}{One-way chain with production}
For species $i=0,1,\ldots$,
\[
 \boxed{N_i'
 =-\lambda_iN_i
 +\lambda_{i-1}N_{i-1}
 +R_i(t)},
 \qquad \lambda_{-1}N_{-1}\equiv0.
\]
Read each equation as
\[
 \text{rate of change}=\text{parent gain}+\text{external gain}-\text{decay loss}.
\]
Initial values $N_i(0)$ complete the system.  Solve the first equation, insert
it into the second, and continue downstream.  Branching ratios multiply the
corresponding parent-gain terms.
\end{fixedbox}

\begin{boardbox}{FNRP extension: secular equilibrium}
For a parent $1$ and daughter $2$ with $N_1(0)=N_{1,0}$, $N_2(0)=0$, no direct
sources, and a 100\% parent-to-daughter branch,
\[
 N_2(t)=\frac{\lambda_1N_{1,0}}{\lambda_2-\lambda_1}
 \left(e^{-\lambda_1t}-e^{-\lambda_2t}\right).
\]
If $\lambda_1\ll\lambda_2$, then after several daughter half-lives but before
the parent changes appreciably,
\[
 A_2(t)\approx\MathBlank[0.66in]{A_1(t)},
 \qquad
 \frac{N_2}{N_1}\approx\MathBlank[0.82in]{\frac{\lambda_1}{\lambda_2}}.
\]
Explain why equal activities do not imply equal inventories.
\RuledLines{1}
\end{boardbox}

\begin{fixedbox}{Steady state versus secular equilibrium}
\textbf{Steady state:} production balances removal for one species, so
$N'=0$ and its inventory is constant.\par
\smallskip
\textbf{Secular equilibrium:} a short-lived daughter tracks a much
longer-lived parent, so their activities are approximately equal while both
may vary slowly.
\end{fixedbox}

\begin{takeawaybox}[Choose the model before solving]
\begin{tabularx}{\linewidth}{@{}>{\bfseries}l X@{}}
no source, one species & $N'=-\lambda N$\\
source, one species & $N'=-\lambda N+R(t)$\\
one-way chain & coupled gain--loss equations\\
long-lived parent & test for secular equilibrium
\end{tabularx}
\end{takeawaybox}

\end{handoutcolumns}
\end{document}
